\chapter{Patterns}

There are three kinds of patterns:
\begin{enumerate}
  \item Theermanam or conclusion,
  \item Korvai or --- google translate gave me a good word --- polynomial,
  \item Korappu or reduction.
\end{enumerate}
Let's look at these one after the other.

\section{Conclusion or Theermanam}

A Theermanam is more like hitting the cymbals at the end of a song. But, this is
structured and mathematically interesting. It is usually either of these,
\begin{itemize}
  \item a combination played thrice;
  \item a pattern that reduces or grows bigger, and then concludes;
  \item a pattern where the second part is the average of the first and third.
\end{itemize}
These definitions are by no means complete or strict. These are merely the most
common ways of playing a Theermanam.

\subsection{Guidelines for creating Theermanams}

The following guidelines might sound confusing. When/if it does, please look at the
examples. I hope they are a lot clearer.
\begin{enumerate}
  \item Decide how many aksharas (claps) the theermanam should last;
  \item Decide the pace --- what nadai (walk) should the piece be in? There might be
    multiple nadai\,s in there. When there is more than one nadai (walk) involved,
    normalize everything to the most used nadai (walk). So, if you have a piece with
    say, chatusra and tisra nadai\,s, normalize the tisram parts to chatusram. This
    is easy. Since the duration of one clap is never altered, use this fact to
    normalize. Three atoms in tisram is exactly the same as four atoms in chatusram.
    So, if you have a part of your theermanam in tisram, and that comes to 18 atoms,
    it would be the same duration as 24 atoms in chatusram.\footnote{This fact can be
    confusing. But, understanding this is essential to ease practice and master
    Carnatic Rhythms.}
  \item Once the above steps are done, you have a specific number of atoms over which
    you can create a theermanam. Write this down.
  \item Come up with a pattern that when done three times over, would cover most of
    the atoms computed above making sure that what is left is an even number and is
    not too large. This would form part of the 'thom'\,s at the end of each pattern.
    The third time, the thom lands on the first clap of the next cycle.\footnote{This
    is called the sama. The thom lands on the sama.}
  \item Divide the balance from the step above by two. Let this be $x$. Insert thoms
    or thams that last for $x$ atoms in between the patterns.
\end{enumerate}

The procedure mentioned above does well for beginners. As you master carnatic rhythm
theory, you would find that steps 1 and 2 are mostly skipped. Intuition takes over.

\subsection{Examples}

\paragraph{The simplest theermanam}
\begin{enumerate}
  \item Lets choose the number of claps --- 4.
  \item Lets choose chatusram --- 4 atoms per clap. The total is 16 atoms. This makes
    the theermanam last for exactly one whole note. Each atom is a
    sixteenth-note.\footnote{I strongly suggest that you don't make a habit of
    converting to western notation every time. This example is simply to ease the
    initial understanding of carnatic terms. I feel that it's easier to use the
    carnatic notation, that is, atoms, claps, et cetera.}
  \item Let the recurring pattern be tha ka dhi na --- 4 atoms. Three times tha ka
    dhi na is 12 atoms.
    \[ 16 \text{ atoms} - 12 \text{ atoms} = 4 \text{ atoms} \]
  \item The subtrahend or balance is 4 atoms. Divide this by 2.
    \[ \frac{4 \text{ atoms}}{2} = 2 \text{ atoms} \]
  \item Every tham or thom has to last for 2 atoms. For convenience --- mostly
    because thom or tham is just one syllable --- let us say that the tham lasts for 1
    atom and the other 1 atom is a kaarvai (a silent syllable).
\end{enumerate}
In the end we obtain,
\begin{center}\small\ttfamily
(Tha ka dhi na | tham) - (tha ka | dhi na tham) - | (tha ka dhi na | tham) \ldots
\end{center}

\paragraph{Extending the simple}
What would happen if we remove the kaarvai\,s from the above 4/4 theermanam? It was 16
atoms long. It would now be 14 atoms long. If we play that at 4 atoms per clap, we
would have 3.5 claps! We have made a theermanam for Misra Chapu. What if we added one
kaarvai after every tham? We would make a theermanam for Sankeerna Chapu. The math is
left as an exercise to the reader.
\begin{center}\small\ttfamily
(Tha ka dhi na | tham) - - (tha | ka dhi na tham) | - - (tha ka | dhi na | tham) \ldots
\end{center}
Note: In the above notation, we have only 18 atoms. Especially, the last clap
(actually half a clap) is only 2 atoms long.

\subsection{Some theermanams from the masters of fusion}

\paragraph{Manirambha conclusion}
The following pattern is from Manirambha (Art Metal). It is a pattern set to Adi Tala
(cycle of 8 claps). To practice and master this one is difficult. However, to follow
and recite it along with the song should be easy. Three types of 'tha dhi ki ta
thom'\,s occur in two different walks each.

\begin{center}
\begin{adjustbox}{max width=\textwidth}
\begin{tabular}{l c}
\toprule
Sol & Claps\\
\midrule
thom tha ka thom tha ka dhi na tham - (10 atoms) & \multirow{3}{*}{$7\frac{1}{2}$ (Chatusram)}\\
thom tha ka thom tha ka dhi na tham - (10 atoms) & \\
thom tha ka thom tha ka dhi na tham - (10 atoms) & \\
\midrule
tha dhi ki ta thom (5 atoms) & \multirow{3}{*}{$3\frac{3}{4}$ (Chatusram)}\\
tha dhi ki ta thom (5 atoms) & \\
tha dhi ki ta thom (5 atoms) & \\
\midrule
tha dhi ki ta thom (5 atoms) & \multirow{3}{*}{$2\frac{1}{2}$ (Tisram)}\\
tha dhi ki ta thom (5 atoms) & \\
tha dhi ki ta thom (5 atoms) & \\
\midrule
tha dhee - ki ta thom (6 atoms) & \multirow{3}{*}{$4\frac{1}{2}$ (Chatusram)}\\
tha dhee - ki ta thom (6 atoms) & \\
tha dhee - ki ta thom (6 atoms) & \\
\midrule
tha dhee - ki ta thom (6 atoms) & \multirow{3}{*}{$3$ (Tisram)}\\
tha dhee - ki ta thom (6 atoms) & \\
tha dhee - ki ta thom (6 atoms) & \\
\midrule
tha - dhee - ki ta thom (7 atoms) & \multirow{3}{*}{$5\frac{1}{4}$ (Chatusram)}\\
tha - dhee - ki ta thom (7 atoms) & \\
tha - dhee - ki ta thom (7 atoms) & \\
\midrule
tha - dhee - ki ta thom (7 atoms) & \multirow{3}{*}{$3\frac{1}{2}$ (Tisram)}\\
tha - dhee - ki ta thom (7 atoms) & \\
tha - dhee - ki ta thom (7 atoms) & \\
\bottomrule
\end{tabular}
\end{adjustbox}
\end{center}

\paragraph{Face to face, Natural Elements, Shakti (2:52 to 3:01)}
The following theermanam is set to Kanda Chapu tala. The walk is tisram (double speed)
--- 6 atoms per clap.

\begin{center}
\begin{adjustbox}{max width=\textwidth}
\begin{tabular}{l c}
\toprule
Sol & Claps\\
\midrule
than gidu thaka thari kita thaka (6 atoms) & 1\\
thom - - (3 atoms = half clap)             & $\frac{1}{2}$\\
\midrule
than gidu thaka thari kita thaka (6 atoms) & \multirow{2}{*}{2}\\
than gidu thaka thari kita thaka (6 atoms) & \\
thom - - (3 atoms = half clap)             & $\frac{1}{2}$\\
\midrule
than gidu thaka thari kita thaka (6 atoms) & \multirow{3}{*}{3}\\
than gidu thaka thari kita thaka (6 atoms) & \\
than gidu thaka thari kita thaka (6 atoms) & \\
thom - - (3 atoms = half clap)             & $\frac{1}{2}$\\
\midrule
tha dhi ki ta thom (5 atoms) & \multirow{3}{*}{$2\frac{1}{2}$}\\
tha dhi ki ta thom (5 atoms) & \\
tha dhi ki ta thom (5 atoms) & \\
\bottomrule
\end{tabular}
\end{adjustbox}
\end{center}

\section{Polynomial or Korvai}

A Korvai is the most interesting of patterns, in my opinion. Especially because the
possibilities are tremendously many. A Korvai has the following chief, but not
strict, characteristics:
\begin{itemize}
  \item has two distinct parts --- poorva anga and uttara anga (Sanskrit) ---
    literally fore and hind parts;
  \item the first part could be a simple or complex combination of fillers;
  \item the second part could be a simple or complex combination of conclusive-fillers.
\end{itemize}

\subsection{Guidelines for creating Korvai\,s}

Since abstract guidelines might be confusing, I will explain the construction of a
korvai with an example:
\begin{itemize}
  \item Choose the tala and nadai (walk) --- keeping it simple, Adi --- 8 clap cycle,
    and Chatusram --- 4 atoms per clap.
  \item Choose a filler --- tha tha kita thaka, thaka thari kita thaka.
  \item Let us create a poorva anga or first part using the filler. The hyphen at the
    end is a kaarvai (swallowed syllable):
\end{itemize}

\begin{center}
\begin{adjustbox}{max width=\textwidth}
\begin{tabular}{l c}
\toprule
Sol & Claps\\
\midrule
tha tha kita thaka, thaka thari kita thaka, thom tha thom - (12 atoms) & \multirow{3}{*}{9}\\
tha tha kita thaka, thaka thari kita thaka, thom tha thom - (12 atoms) & \\
tha tha kita thaka, thaka thari kita thaka, thom tha thom - (12 atoms) & \\
\bottomrule
\end{tabular}
\end{adjustbox}
\end{center}

\begin{itemize}
  \item Clearly, the poorva anga (first part) is 9 claps long. The nearest multiple
    of 8 is 16. The next one is 24. Since 16 seems too short (the second part would
    only have 7 claps then), we choose 24.
    \[ 24 \text{ claps} - 9 \text{ claps} = 15 \text{ claps} \]
  \item As a starting point for the uttara anga (second part) let's pick tha dhee - ki
    ta thom (6 atoms). One combination that would create a uttara anga (second part)
    is a theermanam with conclusive fillers.
\end{itemize}

\begin{center}
\begin{adjustbox}{max width=\textwidth}
\begin{tabular}{l c}
\toprule
Sol & Claps\\
\midrule
tha dhee - ki ta thom (6 atoms) & \multirow{3}{*}{$4\frac{1}{2}$}\\
tha dhee - ki ta thom (6 atoms) & \\
tha dhee - ki ta thom (6 atoms) & \\
Thangu ($3$ atoms or $\frac{3}{4}$ claps) & $\frac{3}{4}$\\
\midrule
tha dhee - ki ta thom (6 atoms) & \multirow{3}{*}{$4\frac{1}{2}$}\\
tha dhee - ki ta thom (6 atoms) & \\
tha dhee - ki ta thom (6 atoms) & \\
Thangu ($3$ atoms or $\frac{3}{4}$ claps) & $\frac{3}{4}$\\
\midrule
tha dhee - ki ta thom (6 atoms) & \multirow{3}{*}{$4\frac{1}{2}$}\\
tha dhee - ki ta thom (6 atoms) & \\
tha dhee - ki ta thom (6 atoms) & \\
\bottomrule
\end{tabular}
\end{adjustbox}
\end{center}

\begin{itemize}
  \item The korvai is already complete. But, as a final touch, we could average out
    the first and last 'tha dhi ki ta thom' triplets. Note that the math is not
    affected by averaging out. Replace the first triple with 'tha dhi ki ta thom' (5
    atoms) and the third triple with 'tha - dhee - ki ta thom' (7 atoms).
\end{itemize}

\paragraph{For the ardent enthusiast}
The above korvai is 24 claps long. It is set to Chatusra nadai (4 atoms per clap). The
total number of atoms is hence,
\[ 24 \text{ claps} \times 4 \text{ atoms/clap} = 96 \text{ atoms} \]
What a miracle! 96 is also 16 times 6. Aha! So could we reinterpret the whole thing as
under?
\[ 16 \text{ claps} \times 6 \text{ atoms/clap} = 96 \text{ atoms} \]
The answer is yes! Playing the whole thing once at 4 atoms per clap, then at 6 atoms
per clap would totally last for 24 + 16 claps. That is 40 claps. You could play that
to Adi Tala (8-clap cycle), or even better, to Kanda chapu (2.5-clap cycle) or Kanda
ekam (5-clap cycle)!!! In essence, we could do the following:
\begin{enumerate}
  \item Play the korvai as such at 4 atoms per clap. This lasts for 24 claps. Fits Adi
    and Rupaka talas.
  \item Play the korvai at 6 atoms per clap. This lasts for 16 claps. Fits Adi tala.
  \item Play the korvai once at 4 atoms per clap, then once at 6 atoms per clap. This
    lasts for 40 claps. Fits Adi and Kanda Chapu/ekam talas.
\end{enumerate}

\subsection{Some korvai\,s from the masters of fusion}

\paragraph{Geometry of Rap, Ragabop Trio (Prasanna, George Brooks, Steve Smith), 2:24 onwards}

\begin{center}
\begin{adjustbox}{max width=\textwidth}
\begin{tabular}{l c}
\toprule
Sol & Claps\\
\midrule
Tha - - - (4 atoms)  & \multirow{4}{*}{4}\\
ka - - - (4 atoms)   & \\
dhee - - - (4 atoms) & \\
mi - - - (4 atoms)   & \\
\midrule
Tha dhee - ki ta thom (6 atoms) & $1\frac{1}{2}$\\
\midrule
Tha - - - (4 atoms)  & \multirow{4}{*}{4}\\
ka - - - (4 atoms)   & \\
dhee - - - (4 atoms) & \\
mi - - - (4 atoms)   & \\
\midrule
Tha dhee - ki ta thom (6 atoms) & \multirow{2}{*}{3}\\
Tha dhee - ki ta thom (6 atoms) & \\
\midrule
Tha - - - (4 atoms)  & \multirow{4}{*}{4}\\
ka - - - (4 atoms)   & \\
dhee - - - (4 atoms) & \\
mi - - - (4 atoms)   & \\
\midrule
Tha dhee - ki ta thom (6 atoms) & \multirow{3}{*}{$4\frac{1}{2}$}\\
Tha dhee - ki ta thom (6 atoms) & \\
Tha dhee - ki ta thom (6 atoms) & \\
Tha - - - (4 atoms or 1 clap)   & 1\\
\midrule
Tha dhee - ki ta thom (6 atoms) & \multirow{3}{*}{$4\frac{1}{2}$}\\
Tha dhee - ki ta thom (6 atoms) & \\
Tha dhee - ki ta thom (6 atoms) & \\
Tha - - - (4 atoms or 1 clap)   & 1\\
\midrule
Tha dhee - ki ta thom (6 atoms) & \multirow{3}{*}{$4\frac{1}{2}$}\\
Tha dhee - ki ta thom (6 atoms) & \\
Tha dhee - ki ta thom (6 atoms) & \\
\bottomrule
\end{tabular}
\end{adjustbox}
\end{center}

\paragraph{La Danse Du Bonheur, A Handful of Beauty, Shakti, 0:10 - 0:24}

\begin{center}
\begin{adjustbox}{max width=\textwidth}
\begin{tabular}{l c}
\toprule
Sol & Claps\\
\midrule
Tha - dhi -, tha ka dhi na, tha ka thom - (12 atoms) & \multirow{2}{*}{6}\\
Tha - dhi -, tha ka dhi na, tha ka thom - (12 atoms) & \\
\midrule
dhi -, tha ka dhi na, tha ka thom - (10 atoms) & \multirow{2}{*}{5}\\
dhi -, tha ka dhi na, tha ka thom - (10 atoms) & \\
\midrule
tha ka dhi na, tha ka thom - (8 atoms) & \multirow{2}{*}{4}\\
tha ka dhi na, tha ka thom - (8 atoms) & \\
\midrule
ju nu, tha ka thom - (6 atoms) & \multirow{2}{*}{3}\\
ju nu, tha ka thom - (6 atoms) & \\
\midrule
tha ka thom - (4 atoms) & \multirow{2}{*}{2}\\
tha ka thom - (4 atoms) & \\
1 atom kaarvai          & \\
\midrule
Tha dhi ki ta thom (5 atoms) & \multirow{3}{*}{$3\frac{3}{4}$}\\
Tha dhi ki ta thom (5 atoms) & \\
Tha dhi ki ta thom (5 atoms) & \\
1 atom kaarvai               & \\
\midrule
Tha dhi ki ta thom (5 atoms) & \multirow{3}{*}{$3\frac{3}{4}$}\\
Tha dhi ki ta thom (5 atoms) & \\
Tha dhi ki ta thom (5 atoms) & \\
1 atom kaarvai               & \\
\midrule
Tha dhi ki ta thom (5 atoms) & \multirow{3}{*}{$3\frac{3}{4}$}\\
Tha dhi ki ta thom (5 atoms) & \\
Tha dhi ki ta thom (5 atoms) & \\
\bottomrule
\end{tabular}
\end{adjustbox}
\end{center}

\paragraph{Carnatic metal, Mattias Eklundh, opening korvai}
This one should serve as a good example to correlate western notation with solkattu
counting since you should have already learnt this one.

\begin{center}
\begin{adjustbox}{max width=\textwidth}
\begin{tabular}{l c}
\toprule
Sol & Claps\\
\midrule
Tha dhem - tha - - thin - (8 atoms)  & \multirow{3}{*}{6}\\
Tha dhem - tha - - thin - (8 atoms)  & \\
Dhin - tha dhin - dhin - - (8 atoms) & \\
\midrule
tha - -   & \multirow{5}{*}{15 atoms or $3\frac{3}{4}$}\\
dhee - -  & \\
ki - -    & \\
ta - -    & \\
thom - -  & \\
\midrule
tha -   & \multirow{5}{*}{10 atoms or $2\frac{1}{2}$}\\
dhee -  & \\
ki -    & \\
ta -    & \\
thom -  & \\
\midrule
tha dhee ki ta thom (5 atoms) & \multirow{3}{*}{$3\frac{3}{4}$}\\
tha dhee ki ta thom (5 atoms) & \\
tha dhee ki ta thom (5 atoms) & \\
1 atom kaarvai                & \\
\midrule
tha dhee ki ta thom (5 atoms) & \multirow{3}{*}{$3\frac{3}{4}$}\\
tha dhee ki ta thom (5 atoms) & \\
tha dhee ki ta thom (5 atoms) & \\
1 atom kaarvai                & \\
\midrule
tha dhee ki ta thom (5 atoms) & \multirow{3}{*}{$3\frac{3}{4}$}\\
tha dhee ki ta thom (5 atoms) & \\
tha dhee ki ta thom (5 atoms) & \\
\bottomrule
\end{tabular}
\end{adjustbox}
\end{center}

\section{Reduction pattern or Korappu}

\subsection{One, one-two, one-two-three, \ldots}

In the last two korvai\,s, did you find something interesting? First, try to separate
the poorva anga (first part) and uttara anga (second part). The poorva anga from the
Shakti korvai and the uttara anga from the Carnatic metal korvai, have something in
common. They are both reduction patterns: while Shakti's reduction pattern is one with
fillers, Carnatic metal's reduction pattern is one with conclusive fillers --- of
course, one is bound to find conclusive fillers in the uttara anga (second part). The
concept is simple, take a filler and slowly shrink it down --- an arithmetic
progression --- the difference between any two parts within the pattern is usually
constant.\footnote{This is also, quite obviously, not a strict rule.} Note that it is
completely valid to move up rather than down. That is, construct a reduction pattern
and then play/recite it in reverse.

Reduction patterns can form parts of a korvai or theermanam. When there is enough
substance in the pattern, and when it is conveniently long, it can go alone in all
it's glory. One such pattern occurs in Big Machine, which you should by now be familiar
with. Since the song itself has solkattu, I will not transcribe it here. It is left as
an exercise to the reader. The filler used in the reduction pattern can be found below
as a hint.\footnote{The filler is: Tha dhi ki ta thom, tha n gu (8 atoms)}

\subsection{Examples of reduction patterns in korvai\,s and theermanams}

\paragraph{Theermanam with a reduction pattern}

\begin{center}
\begin{adjustbox}{max width=\textwidth}
\begin{tabular}{l c}
\toprule
Sol & atoms\\
\midrule
Tha dhi thaka dhina thom tha tho n gu & 9\\
dhi thaka dhina thom tha tho n gu     & 8\\
thaka dhina thom tha tho n gu         & 7\\
thaka thom tha tho n gu               & 6\\
thom tha tho n gu                     & 5\\
tha thom - -                          & 4\\
1 atom kaarvai                        & \\
\midrule
Tha dhi thaka dhina thom tha tho n gu & 9\\
dhi thaka dhina thom tha tho n gu     & 8\\
thaka dhina thom tha tho n gu         & 7\\
thaka thom tha tho n gu               & 6\\
thom tha tho n gu                     & 5\\
tha thom - -                          & 4\\
tha thom - -                          & 4\\
1 atom kaarvai                        & \\
\midrule
Tha dhi thaka dhina thom tha tho n gu & 9\\
dhi thaka dhina thom tha tho n gu     & 8\\
thaka dhina thom tha tho n gu         & 7\\
thaka thom tha tho n gu               & 6\\
thom tha tho n gu                     & 5\\
tha thom - -                          & 4\\
tha thom - -                          & 4\\
tha                                   & 1\\
\bottomrule
\end{tabular}
\end{adjustbox}
\end{center}
thom lands on the first clap of the next avartana (bar).

\paragraph{Korvai with a reduction pattern}

\begin{center}
\begin{adjustbox}{max width=\textwidth}
\begin{tabular}{l c}
\toprule
Sol & atoms\\
\midrule
Tha dhi thaka dhina thom tha thom - & 8\\
Tha dhi thaka dhina thom tha thom - & 8\\
dhi thaka dhina thom tha thom -     & 7\\
dhi thaka dhina thom tha thom -     & 7\\
thaka dhina thom tha thom -         & 6\\
thaka dhina thom tha thom -         & 6\\
thaka thom tha thom -               & 5\\
thaka thom tha thom -               & 5\\
thom tha thom -                     & 4\\
thom tha thom -                     & 4\\
tha thom -                          & 3\\
tha thom -                          & 3\\
thaam -                             & 2\\
thaam -                             & 2\\
\bottomrule
\end{tabular}
\end{adjustbox}
\end{center}

\begin{center}
\begin{adjustbox}{max width=\textwidth}
\begin{tabular}{l c}
\toprule
Sol & Claps\\
\midrule
tha dhee - ki ta thom (6 atoms) & \multirow{3}{*}{$4\frac{1}{2}$}\\
tha dhee - ki ta thom (6 atoms) & \\
tha dhee - ki ta thom (6 atoms) & \\
Tham - (2 atoms)                & \\
\midrule
tha dhee - ki ta thom (6 atoms) & \multirow{3}{*}{$4\frac{1}{2}$}\\
tha dhee - ki ta thom (6 atoms) & \\
tha dhee - ki ta thom (6 atoms) & \\
Tham - (2 atoms)                & \\
\midrule
tha dhee - ki ta thom (6 atoms) & \multirow{3}{*}{$4\frac{1}{2}$}\\
tha dhee - ki ta thom (6 atoms) & \\
tha dhee - ki ta thom (6 atoms) & \\
\bottomrule
\end{tabular}
\end{adjustbox}
\end{center}
thaam lands on the first clap of the next avartana (bar).
